% !TeX root = ../papers/nc-p3.tex
\section{五中心式上の表示と第一長方形の階数十曲面の全正則性}
\label{sec:first-rectangle-edge-regularity}
\begin{theorem}\label{thm:first-rectangle-edge-regularity}
台85のトーラス開集合で $u=1$ または $v=1$ の二曲面は、交線を含め全点が
標準Koszul左右AS正則である。各始点は一般の四中心四次式上で自由階数256である。
$u=-1$ または $v=-1$ の二曲面も全点が標準正則である。
これら四面の指定始点の無印同型はトーラス座標の一致で決まり、周期は
$u=v=\pm1$ で2、残りで4である。
\end{theorem}

\begin{proof}[五つの中心変数によるHilbert上界]
$u=1$, $v\ne1$ として始点 $i$ を固定する。四次中心の五基底を形式変数とする
$S=k[Z_1,\ldots,Z_5]$, $\deg Z_j=4$ を $M=e_iB$ に作用させる。
以下の開集合では全中心商の基底の長さ別個数が
\[
 (q_0,\ldots,q_9)=(1,4,10,20,30,36,34,20,1,0)
\]
となる。基底語を生成元とする自由加群 $L_0$ から $M$ へ次数帰納で全射がある。
$P(T)=(1-T^4)^5/(1-T)^4$ と置き、$q_n=0$ を高次へ延長し
$\beta_n=q_n-[T^n]P(T)$ とする。$\beta_n$ の非零部分は次数8から16まで順に
$1,20,34,36,30,20,10,4,1$ である。核にこれらの個数の $S$ 上独立な斉次関係を取れば、
生成する自由部分加群 $L_1$ に対し $L_0/L_1$ のHilbert級数は $(1-T)^{-4}$ となる。
主成分の幾何学的一般点から各次数に同じ下界があるので $L_0/L_1\to M$ は同型である。
従って $M$ は $S$ 上長さ1の自由分解を持ち、深さと次元がともに4である。
全始点に共通な四つの一般一次結合を次数4の斉次パラメータとすれば、各 $M$ の正則列となる。
その多項式環上の次数付き自由性と階数256は中心商の基底の持上げと級数比較から従う。
補題\ref{lem:parallel-residue-border}の双対完全性の議論により標準Koszul左右AS性を得る。

独立性は五変数の多項式を展開せずに証明できる。
$L_0$ に各矢の作用の $S$ 線型な持上げを選ぶ。具体的には各長さ・各色ごとに、
一文字延長のうち全中心商で消去する列を選ぶ。末尾の二次関係、および長さ4では
中心基底と $Z_j$ 倍の空語との等式を使う。一段短い語の積を先に代入すると、
最高長の係数行列は全中心商での行列であり、選んだ可逆小行列で解ける。
残す列にはその基底語をそのまま用いる。これは $M$ への評価と可換な持上げである。
$Z=c$ で評価した係数だけで同じ消去ができるため、全関係の多項式表示を保存する必要はない。

各基底語での全二次関係、空語での五つの中心等式から出発する。
それらを持上げた右矢と $Z_j$ 倍で閉じれば、常に $L_0\to M$ の核の斉次元を得る。
$V_n$ を内部次数 $n$ で得たこれらの元の $c$ での評価が張る空間とする。
$c_j=1$ となる $j$ を使えば $V_{n-4}\subset V_n$。長さの4による剰余ごとに、
$n\le16$ で $\dim V_n\ge\sum_{a\ge0}\beta_{n-4a}$ を検算する。ただし負の添字では零とする。
前の次数で選んだ独立列を毎回 $\beta_n$ 個だけ延長でき、異なる剰余は異なる基底語を使う。
評価後の独立小行列式が非零なので、元の斉次列は $S$ 上独立である。
\path{computations/quartic_center_presentation.py} が持上げと全次数の独立数を計算する。
\end{proof}

\begin{proof}[二つの曲面の全点]
$U=vz(v-1)(v+1)(v+z)(z-1)(vz+1)$ を逆にする。
表示Aでは位置別文字平行移動 $(0,2,2,0)$ と $c=(1,0,0,0,0)$、
表示Bでは $(2,0,2,0)$ と $c=(0,1,0,0,0)$ を取る。
各表示と始点ごとに、最高長の小行列・関係の独立性に使う小行列の分母とピボットの非単元既約因子を記録する。
Aの四始点での個数は $5,20,28,5$、Bの始点 $0,1,3$ では $7,19,24$ である。
これらは\path{computations/verify_first_rectangle_edge_presentation.py} で再生成できる。

二位置移動は $v\mapsto1/v$ で $z$ を保つ。その引戻しを上線で表す。
始点0は $A_0,B_0,\overline{A_2}$ で覆われる。実際、最初の二リストの因子の35対ごとに
二因子と $1-TU$ のイデアルに第三リストの因子積を加えると単位イデアルである。
始点1は $v+2z-1\ne0$ で $\overline{A_3},B_1,A_1,\overline{B_3}$ で覆われる。
最初の二リストの95対に対する同じ計算で、今度は $U(v+2z-1)$ を逆にする。
残る直線では $z=(1-v)/2$ を $\overline{A_3}$ の各因子に代入すると、単元を除く零点は
\[
 v^3-4v^2+3v-4,\quad v^2-2,\quad v^2+v+1,\quad 2v^2-3v+2
\]
の四数体に限られる。順に $(p,v)=(1013,412),(1009,570),(1009,634),(1009,274)$ は
単純な剰余根である。各点で四中心式の商を全始点の長さ13まで検算し、前節と同じ持上げで
標数零の全共役点に自由性と左右正則性を得る。
\path{computations/verify_first_rectangle_edge_regularity.py} が35対・95対および四数体を確かめる。
二位置移動で始点2,3も従い、一位置移動で $v=1$ の面へ移る。
交線と負側はそれぞれ既知の垂直成分の定理と四元群二コサイクル捻りによる。
同型と周期は定理\ref{thm:first-rectangle-horizontal-frames}による。
\end{proof}
